\documentclass[11pt,a4paper]{article}

\usepackage[utf8]{inputenc}
\usepackage[T1]{fontenc}
\usepackage[spanish]{babel}
\usepackage{amsmath,amssymb,mathtools}
\usepackage{geometry}
\usepackage{enumitem}
\usepackage{array}
\usepackage{lmodern}
\usepackage{microtype}
\usepackage{xcolor}

\geometry{margin=2.4cm}
\setlength{\parindent}{0pt}
\setlength{\parskip}{0.5em}

\newcommand{\R}{\mathbb{R}}
\newcommand{\K}{\mathbb{K}}
\newcommand{\Pn}{\mathcal{P}}
\newcommand{\Span}{\operatorname{L}}
\newcommand{\rg}{\operatorname{rg}}
\newcommand{\dimn}{\operatorname{dim}}

\begin{document}

\begin{center}
	{\LARGE\bfseries Álgebra Lineal}\\[0.4em]
	{\Large Tema 1: Espacios Vectoriales --- Ejercicios}
\end{center}

\section{Ejercicio 1}
Razonar si el siguiente conjunto constituye un subespacio vectorial. En caso afirmativo, encontrar su expresión como subespacio engendrado por un sistema de vectores:

\[
S=
\left\{
(x,y,z)\in\R^3:
2x-y+3z=0
\right\}
\]

\section{Ejercicio 2}
Razonar si el siguiente conjunto constituye un subespacio vectorial. En caso afirmativo, encontrar su expresión como subespacio engendrado por un sistema de vectores:

\[
S=
\left\{
(x,y)\in\R^2:
xy=0
\right\}
\]

\section{Ejercicio 3}
Razonar si el siguiente conjunto constituye un subespacio vectorial:

\[
S=
\left\{
A=
\begin{pmatrix}
a&b\\
c&d
\end{pmatrix}
\in M_{2\times2}(\R):
\det A=0
\right\}
\]

\section{Ejercicio 4}
Determinar el valor de $x$ para que:

\[
(1,x,5)\in
\Span\left\{
(1,2,3),(1,1,1)
\right\}
\]

\section{Ejercicio 5}
Determinar si los conjuntos:

\[
A=
\left\{
p(x)\in\Pn_2(x):
p(0)=0,\;
p'(0)=0
\right\}
\]

y:

\[
B=
\left\{
p(x)\in\Pn_2(x):
p(0)=0,\;
p'(0)=1
\right\}
\]

son subespacios de $\Pn_2(x)$.

\section{Ejercicio 6}
Obtener las ecuaciones paramétricas del subespacio:

\[
M=
\left\{
(x_1,x_2,x_3,x_4)\in\R^4:
\begin{array}{rcl}
x_1+5x_2-2x_3-x_4&=&0\\
2x_1+x_2+2x_3+x_4&=&0
\end{array}
\right\}
\]

\section{Ejercicio 7}
Dados los subespacios:

\[
S=
\left\{
(x_1,x_2,x_3,x_4)\in\R^4:
x_1-x_2=0
\right\}
\]

\[
T=
\Span\left\{
(1,1,2,1),
(2,3,-1,1)
\right\}
\]

obtener bases de $S$, $T$, $S\cap T$ y $S+T$.

\section{Ejercicio 8}
Analizar para qué valores de $a$ las siguientes matrices son linealmente independientes:

\[
A_1=
\begin{pmatrix}
0&a&0\\
0&1&0
\end{pmatrix},
\qquad
A_2=
\begin{pmatrix}
0&1&0\\
0&a&0
\end{pmatrix},
\qquad
A_3=
\begin{pmatrix}
0&1&0\\
0&a&a
\end{pmatrix}
\]

\section{Ejercicio 9}
Analizar para qué valores de $a$ los siguientes polinomios son linealmente independientes:

\[
x^2+3x+1,
\qquad
2-x,
\qquad
1+ax+x^2
\]

\section{Ejercicio 10}
Se considera la matriz:

\[
A=
\begin{pmatrix}
1&-1&1\\
1&-1&1
\end{pmatrix}
\]

Hallar la dimensión y una base del subespacio:

\[
U=
\left\{
X\in M_{2\times2}(\R):
XA=0
\right\}
\]

\section{Ejercicio 11}
En el espacio vectorial $\R^4$, calcular:

\begin{enumerate}[label=\alph*)]
	\item una base que contenga al vector $(1,2,1,1)$;
	\item una base que contenga a los vectores $(1,1,0,0)$, $(0,0,2,2)$ y $(0,3,3,0)$.
\end{enumerate}

\section{Ejercicio 12}
Dadas:

\[
B=\{e_1,e_2,e_3\}
\]

\[
B'=
\left\{
2e_1+3e_2,\;
e_1+e_3,\;
-e_2+e_3
\right\}
\]

se pide:

\begin{enumerate}[label=\alph*)]
	\item Si $[u]_B=(1,2,3)$, hallar $[u]_{B'}$.
	\item Si $[u]_{B'}=(-2,1,0)$, hallar $[u]_B$.
\end{enumerate}

\section{Ejercicio 13}
Sea $\Pn_3(x)$ con bases:

\[
B=(1,x,x^2,x^3)
\]

\[
B'=
(1+x,\;x+x^2,\;x^2-x^3,\;1+2x^3)
\]

\[
B''=
(1,\;1+x,\;1+x+x^2,\;1+x+x^2+x^3)
\]

Se pide:

\begin{enumerate}[label=\alph*)]
	\item Hallar la matriz de cambio de base de $B$ a $B'$.
	\item Hallar la matriz de cambio de base de $B$ a $B''$.
	\item Hallar la matriz de cambio de base de $B'$ a $B''$.
	\item Para $p(x)=x^3-2x$, hallar sus coordenadas respecto de $B$, $B'$ y $B''$.
\end{enumerate}

\section{Ejercicio 14}
En $\Pn_2(x)$ se consideran las bases:

\[
B_1=
\{1+x+x^2,\;x+2x^2,\;1+x\}
\]

\[
B_2=
\{1+2x+3x^2,\;
\alpha+(\alpha-1)x-2x^2,\;
2+2x\}
\]

Se pide:

\begin{enumerate}[label=\alph*)]
	\item Calcular $\alpha$ para que el polinomio de coordenadas $(1,1,0)$ en $B_2$ tenga coordenadas $(1,0,3)$ en $B_1$.
	\item Para dicho valor de $\alpha$, determinar el conjunto $W$ de polinomios que tienen las mismas coordenadas en $B_1$ y $B_2$. Determinar si $W$ es subespacio vectorial de $\Pn_2(x)$.
	\item En caso afirmativo, obtener unas ecuaciones implícitas en base $B_1$ de un subespacio suplementario de $W$ en $\Pn_2(x)$.
\end{enumerate}

\section{Ejercicio 15}
En $\Pn_3(x)$ se consideran:

\[
S_1=
\left\{
p(x)\in\Pn_3(x):
p(0)=0
\text{ y las tangentes a }p(x)
\text{ en }x=1\text{ y }x=-1
\text{ son paralelas}
\right\}
\]

\[
S_2=
\left\{
p(x)\in\Pn_3(x):
p(2)=0
\right\}
\]

Se pide:

\begin{enumerate}[label=\alph*)]
	\item Calcular la dimensión y una base de $S_1$ y $S_2$.
	\item Calcular $S_1\cap S_2$, una base del mismo y razonar si $S_1+S_2$ es suma directa.
	\item Si $p(x)$ es un vector de $S_1\cap S_2$, calcular sus coordenadas en:
	\[
	B=
	\{1+x,\;x+x^2,\;x^2+x^3,\;x^3\}
	\]
	\item Comprobar mediante la matriz de cambio de base apropiada que las coordenadas obtenidas corresponden al mismo vector en la base canónica.
\end{enumerate}

\section{Ejercicio 16}
Sea $\Pn_2(x)$ y sea $p(x)$ un polinomio de grado exactamente $2$.

\begin{enumerate}[label=\alph*)]
	\item Demostrar que:
	\[
	B=\{p,p',p''\}
	\]
	y:
	\[
	B'=\{p,\;p+p',\;p'+p''\}
	\]
	son bases de $\Pn_2(x)$.
	\item Analizar si:
	\[
	S=
	\{x(x-a):a\in\R\}
	\]
	es un subespacio vectorial de $\Pn_2(x)$.
	\item Si $p(x)\in S$ y tiene raíz $-1$, obtener las coordenadas de:
	\[
	q(x)=x^2+x+2
	\]
	en la base $B'$.
	\item Comprobar mediante una matriz de cambio de base que el vector obtenido es precisamente $q(x)$.
\end{enumerate}

\section{Ejercicio 17}
Se consideran los subespacios de $\R^4$:

\[
V=
\Span\{(1,2,3,4),(-1,0,1,-1)\}
\]

\[
W=
\left\{
(x,y,z,t)\in\R^4:
2x+5y-z-t=0
\right\}
\]

Se pide:

\begin{enumerate}[label=\alph*)]
	\item Obtener una base de $W$.
	\item Obtener las ecuaciones paramétricas, implícitas y una base de $V\cap W$.
	\item Razonar si $V+W$ es suma directa.
	\item Hallar las coordenadas de $(-1,1,1,2)$ respecto de la base de $V$ dada en el enunciado.
\end{enumerate}

\section{Ejercicio 18}
En $\Pn_3(x)$ se consideran:

\[
V=
\left\{
p(x)\in\Pn_3(x):
p'(x)\in
\Span\{1+x^2,x^3\}
\right\}
\]

\[
W=
\left\{
p(x)\in\Pn_3(x):
p'(x)=p''(x)
\right\}
\]

Se pide:

\begin{enumerate}[label=\alph*)]
	\item Calcular una base y las ecuaciones implícitas de $V$ y $W$ en la base canónica de $\Pn_3(x)$.
	\item Calcular las ecuaciones paramétricas y una base de $V\cap W$.
	\item Determinar si $p(x)=1$ pertenece a $V$ y a $V\cap W$. En caso afirmativo, calcular sus coordenadas en las bases obtenidas.
\end{enumerate}

\section{Ejercicio 19}
Sea $\Pn_n(x)$ el espacio vectorial de polinomios de grado menor o igual que $n$.

\begin{enumerate}[label=\alph*)]
	\item Demostrar que $x^n$ y sus $n$ primeras derivadas conforman una base de $\Pn_n(x)$.
	\item Estudiar si:
	\[
	1+3x+5x^2,\qquad
	-1+2x^2,\qquad
	3+3x+x^2
	\]
	son linealmente independientes.
	\item Si:
	\[
	V=
	\Span\{1+x^2,1-x^2\}
	\]
	determinar si $p(x)=1+5x^2$ y $r(x)=1+x$ pertenecen a $V$.
	\item Si:
	\[
	W=
	\Span\{1+3x+5x^2,\;-1+2x^2,\;3+3x+x^2\}
	\]
	calcular $V\cap W$ y $V+W$.
\end{enumerate}

\section{Ejercicio 20}
Sean $E=F=\R^2$. Calcular una base del espacio vectorial producto $E\times F$ asociada a la base canónica de $\R^2$.

\section{Ejercicio 21}
Sea $E$ un espacio vectorial de dimensión $3$ y:

\[
B=\{e_1,e_2,e_3\}
\]

una base de $E$. Sea:

\[
V=
\Span\{u_1,u_2\}
\]

con:

\[
u_1=e_1-e_2,
\qquad
u_2=e_1+e_2
\]

Calcular una base del espacio vectorial cociente $E/V$.

\section{Ejercicio 22}
Obtener una base del espacio vectorial cociente $\R^4/V$, siendo:

\[
V=
\Span\left\{
(1,0,1,1),
(1,2,1,1),
(2,2,2,2)
\right\}
\]

\end{document}
